% Single-file paper. Edge cases showcased:
%   - section-based numbering (theorem 1.1, 2.1, ...)
%   - a lemma sharing the theorem counter (1.2)
%   - an unnumbered environment (\newtheorem*)
%   - a proof immediately following its theorem
%   - user macros (\newcommand, with and without arguments) to be expanded
\documentclass{article}
\usepackage{amsmath, amsthm, amssymb}

\newtheorem{theorem}{Theorem}[section]
\newtheorem{lemma}[theorem]{Lemma}
\newtheorem{definition}[theorem]{Definition}
\newtheorem*{remark}{Remark}

\newcommand{\R}{\mathbb{R}}
\newcommand{\norm}[1]{\left\lVert #1 \right\rVert}

\begin{document}

\section{Foundations}

\begin{definition}\label{def:norm}
A \emph{norm} on $\R^n$ is a function $\norm{\cdot} : \R^n \to \R$ that is
positive-definite, homogeneous, and satisfies the triangle inequality.
\end{definition}

\begin{theorem}[Nonnegativity]\label{thm:main}
For every $x \in \R^n$ we have $\norm{x} \geq 0$, with equality iff $x = 0$.
\end{theorem}

\begin{proof}
Immediate from positive-definiteness in Definition~\ref{def:norm}.
\end{proof}

\begin{lemma}\label{lem:aux}
There exists $y \in \R^n$ with $\norm{y} > 0$ whenever $n \geq 1$.
\end{lemma}

\section{Remarks}

\begin{theorem}
A second numbered theorem, appearing in the second section.
\end{theorem}

\begin{remark}
This unnumbered remark should carry no reference number.
\end{remark}

\end{document}
